%=====================================================================
%  LaTeX‑dokument – norsk tekst om Hill‑kryptering (klar for Overleaf)
%=====================================================================
\documentclass[12pt,a4paper]{article}

%-------------------------- Pakker ------------------------------------
\usepackage[utf8]{inputenc}          % tegnkoding
\usepackage[T1]{fontenc}             % korrekt tegnbehandling
\usepackage[norsk]{babel}            % språk
\usepackage{amsmath,amssymb,amsthm}  % matematiske omgivelser
\usepackage{graphicx}                % figurer (ikke brukt her)
\usepackage{float}                   % presis plassering av figurer/tabeller
\usepackage{booktabs}                % pene tabeller
\usepackage{hyperref}                % lenker
\usepackage{geometry}
\geometry{margin=2.5cm}

%-------------------------- Dokumentinformasjon -----------------------
\title{Prosjekt 1: Hill-siffer}
\author{ Maria Brizida Inacio Dos Santos Carvalho, Stian Frøland \\ Kristina Løvbrekke, Linnea Nordstrand, Karoline Låge Sivertsen }
\date{September 2026}
\setlength{\parindent}{0pt}


%=====================================================================
\begin{document}
\maketitle
\tableofcontents
\newpage

%=====================================================================
\section{Innledning}
Kryptering er en metode for å omgjøre data, i disse eksemplene ren tekst,
til chiffertekst ved bruk av en krypteringsnøkkel. Dekryptering gjør
chiffertekst om til den originale rene teksten. Formålet med kryptering og
dekryptering er å dele opp data på en sikker måte, ved at den kun er lesbar
for de med krypteringsnøkkel, og fremstår uforståelig for alle andre.
Dette er kritisk for å kunne sende over viktig informasjon, uten å få den
på avveie.

%=====================================================================
\section{Alfabet og tall}
Man kan omgjøre det norske alfabetet til sifre $0$–$28$ der
$
A=0,\; B=1,\; C=2,\;\dots ,\; Å=28 .
$
Vi tar for oss vektoren
$
\mathbf{H} = (H,\;E,\;I,\;V,\;E,\;R,\;D,\;E,\;N).
$

Siden hver av bokstavene har et tilsvarende siffer kan vi omgjøre vektoren
med bokstaver til en vektor med tall:
$
H=7,\;E=4,\;I=8,\;V=21,\;R=17,\;D=3,\;N=13 .
$

Dermed blir tallvektoren
$
\mathbf{v} = (7,\,4,\,8,\,21,\,4,\,17,\,3,\,4,\,13).
$

\noindent Oversikten $A=0$, $B=1$, osv. fungerer som krypteringsnøkkelen i dette
tilfellet.

%=====================================================================
\section{Modulo‑operasjon}
\texttt{Modulo}(x) betyr at resultatet ligger i mengden $\{0,1,\dots ,x-1\}$;
svaret kan altså ikke være $x$ eller høyere. I eksempelet der modulo er $29$,
må resultatet være under $29$.

Vi kommer fram til at svaret er $85$, men siden $85\ \text{mod }29$ ikke er
mulig direkte, ser vi på hvor mye $85$ overskrider $29$:
$
85 = 2\cdot 29 + 27 \quad\Longrightarrow\quad 85\ (\text{mod }29)=27 .
$

%=====================================================================
\section{Hill‑kryptering}
\subsection{Generell fremgangsmåte}
Måten å kryptere en streng på via Hill‑siffer er å først velge en nøkkel
$N$. Nøkkelen er en matrise, for eksempel en $(3\times3)$‑matrise.
Deretter trenger du meldingen $M$, som du vil kryptere. Meldingen blir delt opp
i blokker som passer nøkkelmatrisen slik at den kan ganges med den.

Bokstavene i matrisen gjøres om til tall ($A=0$, $B=1$, osv.). Deretter
ganges nøkkelmatrisen $N$ med meldingsmatrisen $M$. Hvert element i
produktmatrisen $NM$ regnes modulo $p$, der $p$ er alfabetstørrelsen (for
det norske alfabetet $p=29$). Til slutt gjøres tallene om tilbake til bokstaver,
og man får den krypterte meldingen.

%=====================================================================
\section{Oppgave}
Vårt ord på $4$ bokstaver er \texttt{SMIL}. Ordets vektor er
$
\mathbf{v}= (18,\ 12,\ 8,\ 11).
$

Vi bruker modulo $29$ med matrisen
$
A=
\begin{bmatrix}
2 & 4 & 0 & 8 \\
0 & 3 & 2 & 5 \\
0 & 0 & 9 & 8 \\
0 & 0 & 6 & 3
\end{bmatrix}.
$

Produktet $A\mathbf{v}$ beregnes rad for rad:

$
\begin{aligned}
\text{Av}_{1}&=2\cdot18+4\cdot12+0\cdot8+8\cdot11
               =172\pmod{29}=27,\\[4pt]
\text{Av}_{2}&=0\cdot18+3\cdot12+2\cdot8+5\cdot11
               =107\pmod{29}=20,\\[4pt]
\text{Av}_{3}&=0\cdot18+0\cdot12+9\cdot8+8\cdot11
               =160\pmod{29}=15,\\[4pt]
\text{Av}_{4}&=0\cdot18+0\cdot12+6\cdot8+3\cdot11
               =177\pmod{29}=3.
\end{aligned}
$

Dermed får vi vektoren
$
A\mathbf{v}= \begin{bmatrix}27\\20\\15\\3\end{bmatrix},
$
som kan skrives om til strengen \texttt{ØUPD}.

%=====================================================================
\subsection{4b) Inversibilitet}
Nøkkel‑matrisen må være inverterbar for å kunne dekryptere meldingen.
Uten en inversmatrise er det ingen måte å gå tilbake fra chiffer til
klartekst.  

Det norske alfabetet har $29$ bokstaver; når man regner modulo $29$ har hver
bokstav et tilsvarende siffer. Alle meldinger på norsk kan krypteres ved å
bruke modulo $29$, forutsatt at matrisen $A$ er inverterbar.

Matrisen $A$ er inverterbar dersom
$
\det(A)\not\equiv 0 \pmod{29}.
$
Hvis $\det(A)\equiv0\ (\text{mod }29)$ finnes ingen invers, og kryptering med
den matrisen kan ikke dekrypteres \cite{minini2026}.

\[ A \text{ er inverterbar dersom }\det(A)\neq 0 \] \[ A= \begin{pmatrix} 2 & 4 & 0 & 8\\ 0 & 3 & 2 & 5\\ 0 & 0 & 9 & 8\\ 0 & 0 & 6 & 3 \end{pmatrix} \] \[ M_{11}= \begin{vmatrix} 3 & 2 & 5\\ 0 & 9 & 8\\ 0 & 6 & 3 \end{vmatrix} \qquad M_{21}= \begin{vmatrix} 4 & 0 & 8\\ 0 & 9 & 8\\ 0 & 6 & 3 \end{vmatrix} \] \[ M_{31}= \begin{vmatrix} 4 & 0 & 8\\ 3 & 2 & 5\\ 0 & 6 & 3 \end{vmatrix} \qquad M_{41}= \begin{vmatrix} 4 & 0 & 8\\ 3 & 2 & 5\\ 0 & 9 & 8 \end{vmatrix} \] \[ \det(A) = 2\cdot M_{11} + 0\cdot M_{21} + 0\cdot M_{31} + 0\cdot M_{41} \] \[ = 2\cdot \left( 3\cdot(27-48) - 2\cdot(0) + 5\cdot(0) \right) \] \[ = 2\cdot(-21\cdot3) = 2\cdot(-63) = -126 \] \[ -126 \equiv -10 \pmod{29} \] \[ -10 \neq 0 \quad\Rightarrow\quad \text{inverterbar mod }29 \]

%=====================================================================
\subsection{4c) Kontroll av invers}
For å sjekke om matrisen er invers, kan man:

\begin{enumerate}
\item Rad‑redusere den sammen med identitetsmatrisen.
\item Multiplikere $A$ med en foreslått invers $A^{-1}$ og se om resultatet er
      identitetsmatrisen (modulo $29$).
\end{enumerate}

I oppgaven ble følgende rad‑reduksjon utført (forkortet for lesbarhet):

$
\begin{bmatrix}
2&4&0&8&\;1&0&0&0\\
0&3&2&5&\;0&1&0&0\\
0&0&9&8&\;0&0&1&0\\
0&0&6&3&\;0&0&0&1
\end{bmatrix}
\;\longrightarrow\;
\begin{bmatrix}
2&4&0&8&\;1&0&0&0\\
0&3&2&5&\;0&1&0&0\\
0&0&1&\tfrac{8}{9}&\;0&0&\tfrac{1}{9}&0\\
0&0&0&1&\;0&0&\tfrac{2}{7}&-\tfrac{3}{7}
\end{bmatrix}.
$

Den nederste raden viser pivotelementet $1$ og koeffisientene
$\frac{2}{7}$ og $-\frac{3}{7}$, som i modulo $29$ tilsvarer $21$ og $12$.
Siden dette ikke stemmer med de faktiske elementene i $A^{-1}$, er
$A^{-1}$ **ikke** inversen av $A$.

En direkte kontroll ved multiplikasjon gir:

$
AA^{-1}=
\begin{bmatrix}
15&9&1&10\\
0&10&1&0\\
0&0&4&28\\
0&0&21&12
\end{bmatrix}
\begin{bmatrix}
2&4&0&8\\
0&3&2&5\\
0&0&9&8\\
0&0&6&3
\end{bmatrix}
=
\begin{bmatrix}
30&147&\dots\\
\vdots & \vdots & \ddots
\end{bmatrix}
\pmod{29}
=
\begin{bmatrix}
1&2&\dots\\
\vdots&\vdots&\ddots
\end{bmatrix},
$
hvor elementet $(1,2)$ er $2$ i stedet for $0$, så produktet er **ikke**
identitetsmatrisen. Dermed er den foreslåtte matrisen ikke den korrekte
inversen.

\subsubsection*{Dekryptering}
For å dekryptere en kryptert melding trenger man:
\begin{enumerate}
\item Den krypterte meldingen (omgjort til tall‑matrise).
\item Inversen av nøkkelmatrisen $A^{-1}$ (modulo $29$).
\end{enumerate}
Man ganger $A^{-1}$ med den krypterte meldingsmatrisen, tar resultatet
modulo $29$ og oversetter tallene tilbake til bokstaver.  

Den korrekte inversen av $A$ (modulo $29$) finnes ved å bruke en algoritme for
modulær matrix‑invers, men den er **ikke** den som er oppgitt i oppgaven
over.

%=====================================================================
\section*{Bibliografi}
\begin{thebibliography}{9}
\bibitem{minini2026}
Minini, A. (2026, 16.~september). \emph{Modular Inverse of a Matrix}. Hentet fra
Andrea Minini: \url{https://www.andreaminini.net/math/modular-inverse-of-a-matrix}.

\bibitem{wikipedia}
Hill cipher. \emph{Wikipedia}. \url{https://en.wikipedia.org/wiki/Hill_cipher}.

\bibitem{geeksforgeeks}
Hill Cipher – Matrix Multiplicative Inverse Example.
\url{https://www.geeksforgeeks.org/dsa/hill-cipher/}.

\bibitem{cybrary}
Hill Cipher – 3×3 Matrix Multiplicative Inverse.
\url{https://www.cybrary.it/blog/learn-hill-cipher-3x3-matrix-multiplicative-inverse-example}.
\end{thebibliography}

\end{document}
